% Purpose: Formal FDSN citations, persistent identifiers (DOIs), and data availability statements for all 18 seismic networks utilized in this study.
% Status: Draft; author and venue review pending.
% Source-of-truth: FDSN registry, EIDA portal, OGS metadata archives.

\label{chap:data_resources}

This chapter records the intended data-provenance and software-availability statement for this dissertation. The network identifiers and links below require human confirmation against the FDSN registry and the access date used for the submitted thesis. They are not validated by the source checkout, and restricted waveforms or institutional catalogs are not redistributed with it.

\section*{Seismic Waveform and Station Metadata Networks}
Continuous seismic waveform time series and station response inventories (StationXML) were retrieved from distributed European \ac{FDSN} data centers (including OGS, \ac{INGV}, GeoSphere Austria, the Swiss Seismological Service \ac{SED}, the Slovenian Environment Agency \ac{ARSO}, and the \ac{ORFEUS} \ac{EIDA} infrastructure). 

A total of \textbf{18 distinct regional, national, and international seismic networks} contributed observational data to the 266-station monitoring array analyzed in this thesis. Table~\ref{tab:network_inventory} provides a systematic synthesis of these network assets, followed by formal bibliographic and data repository citations.

\begin{table}[htbp]
  \centering
  \small
  \caption{Inventory of participating seismic networks, operating institutions, and persistent identifiers.}
  \label{tab:network_inventory}
  \begin{threeparttable}
    \begin{tabularx}{\linewidth}{l l l c X}
      \toprule
      \textbf{Code} & \textbf{Network Name} & \textbf{Operating Institution} & \textbf{Period} & \textbf{Persistent Identifier / DOI} \\
      \midrule
      \textbf{4P} & AdriaArray Temporary Network (N/S) & AdriaArray Working Group & 2022--2026 & \url{https://www.fdsn.org/networks/detail/4P_2022/} \\
      \textbf{2Y} & AdriaArray Temporary Network (NE Italy) & OGS & 2022--2026 & \url{https://doi.org/10.7914/1p36-6t87} \\
      \textbf{CH} & National Seismic Networks of Switzerland & SED at ETH Z\"urich & 1983--pres. & \url{https://doi.org/10.12686/sed/networks/ch} \\
      \textbf{EV} & Collalto Seismic Network (CSN) & OGS & 2012--pres. & \url{https://doi.org/10.7914/SN/EV} \\
      \textbf{GU} & Regional Seismic Network of NW Italy & University of Genoa & 1967--pres. & \url{https://doi.org/10.7914/SN/GU} \\
      \textbf{HA} & Interreg Italia-Austria Network & OGS / ZAMG & 2005--pres. & FDSN temporary / project archive \\
      \textbf{IT} & Italian Strong Motion Network (RAN) & Civil Protection Department & 1972--pres. & \url{https://doi.org/10.7914/SN/IT} \\
      \textbf{IV} & Rete Sismica Nazionale (RSN) & INGV & 2005--pres. & \url{https://doi.org/10.13127/sd/x0fxnh7qfy} \\
      \textbf{MN} & Mediterranean Network (MedNet) & INGV / Partners & 1990--pres. & \url{https://doi.org/10.13127/sd/fbbbtdtd6q} \\
      \textbf{NI} & North-East Italy Broadband Network & OGS / Univ.\ of Trieste & 2002--pres. & \url{https://doi.org/10.7914/SN/NI} \\
      \textbf{OE} & Austrian Seismic Network & GeoSphere Austria (ZAMG) & 1987--pres. & \url{https://doi.org/10.7914/SN/OE} \\
      \textbf{OX} & North-East Italy Seismic Network & OGS & 2016--pres. & \url{https://doi.org/10.7914/SN/OX} \\
      \textbf{RF} & FVG Accelerometric Network (RAF) & University of Trieste & 1993--pres. & \url{https://doi.org/10.7914/SN/RF} \\
      \textbf{S}  & Seismo-at-School Program & SED at ETH Z\"urich & 2008--pres. & \url{https://doi.org/10.12686/SED/NETWORKS/S} \\
      \textbf{SI} & Province S\"udtirol Network & ZAMG / Prov.\ Bozen & 2006--pres. & \url{https://www.fdsn.org/networks/detail/SI/} \\
      \textbf{SL} & Seismic Network of the Republic of Slovenia & Slovenian Env.\ Agency (ARSO) & 1990--pres. & \url{https://doi.org/10.7914/SN/SL} \\
      \textbf{ST} & Trentino Seismic Network & Prov.\ Autonoma di Trento & 1981--pres. & \url{https://doi.org/10.7914/SN/ST} \\
      \textbf{ZO} & Garda Multiparametric Network (PDnet) & INGV & 2021--pres. & \url{https://doi.org/10.13127/SD/YHCFOMCBO_} \\
      \bottomrule
    \end{tabularx}
    \begin{tablenotes}[flushleft]
      \footnotesize
      \item Target \& Scope: Permanent and temporary seismic monitoring networks utilized in the 266-station case study array.
      \item What the reader should notice: The dataset integrates 18 distinct network operators across Italy, Austria, Slovenia, and Switzerland into a unified transboundary observation aperture.
      \item Evidence Anchor: Formally registered FDSN network codes and International DOI Foundation records.
    \end{tablenotes}
  \end{threeparttable}
\end{table}

\subsection*{Formal Network Citations}

\begin{enumerate}[label={[\textbf{\arabic*}]},leftmargin=*]
  \item \textbf{[4P] AdriaArray Working Group (2022)}. 
  \textit{AdriaArray Temporary Network: Italy -- north, south (2022--2026)}. 
  International Federation of Digital Seismograph Networks. 
  \url{https://www.fdsn.org/networks/detail/4P_2022/}

  \item \textbf{[2Y] Pesaresi, D., \& Rossi, G. (2022)}. 
  \textit{AdriaArray Temporary Network: Italy -- northeast (OGS-AdA)} [Data set]. 
  International Federation of Digital Seismograph Networks. 
  \url{https://doi.org/10.7914/1p36-6t87}

  \item \textbf{[CH] Swiss Seismological Service (SED) at ETH Z\"urich (1983)}. 
  \textit{National Seismic Networks of Switzerland} [Data set]. 
  ETH Z\"urich. 
  \url{https://doi.org/10.12686/sed/networks/ch}

  \item \textbf{[EV] Priolo, E., Romanelli, M., Plasencia Linares, M.~P., Garbin, M., Peruzza, L., Romano, M.~A., Marotta, P., Bernardi, P., Moratto, L., Zuliani, D., \& Fabris, P. (2015)}. 
  \textit{Seismic Monitoring of an Underground Natural Gas Storage Facility: The Collalto Seismic Network}. 
  \textit{Seismological Research Letters}, 86(1), 109--123. 
  \url{https://doi.org/10.1785/0220140087}; 
  \textit{Collalto Seismic Network (CSN)} [Data set]. International Federation of Digital Seismograph Networks. 
  \url{https://doi.org/10.7914/SN/EV}

  \item \textbf{[GU] University of Genoa (1967)}. 
  \textit{Regional Seismic Network of North Western Italy (RSNI)} [Data set]. 
  International Federation of Digital Seismograph Networks. 
  \url{https://doi.org/10.7914/SN/GU}

  \item \textbf{[HA] OGS and ZAMG (2005)}. 
  \textit{Interreg Italia-Austria: Historical And Recent Earthquakes in Italy and Austria (HA)}. 
  Cross-border seismological data exchange platform. 
  OGS Seismological Research Centre, Udine/Trieste.

  \item \textbf{[IT] Presidency of Council of Ministers -- Civil Protection Department (1972)}. 
  \textit{Italian Strong Motion Network (Rete Accelerometrica Nazionale -- RAN)} [Data set]. 
  International Federation of Digital Seismograph Networks. 
  \url{https://doi.org/10.7914/SN/IT}

  \item \textbf{[IV] Istituto Nazionale di Geofisica e Vulcanologia (INGV) (2005)}. 
  \textit{Rete Sismica Nazionale (RSN)} [Data set]. 
  Istituto Nazionale di Geofisica e Vulcanologia (INGV). 
  \url{https://doi.org/10.13127/sd/x0fxnh7qfy}

  \item \textbf{[MN] MedNet Project Partner Institutions (1990)}. 
  \textit{Mediterranean Very Broadband Seismographic Network (MedNet)} [Data set]. 
  Istituto Nazionale di Geofisica e Vulcanologia (INGV). 
  \url{https://doi.org/10.13127/sd/fbbbtdtd6q}

  \item \textbf{[NI] OGS (Istituto Nazionale di Oceanografia e di Geofisica Sperimentale) and University of Trieste (2002)}. 
  \textit{North-East Italy Broadband Network} [Data set]. 
  International Federation of Digital Seismograph Networks. 
  \url{https://doi.org/10.7914/SN/NI}

  \item \textbf{[OE] ZAMG -- Zentralanstalt f\"ur Meteorologie und Geodynamik (1987)}. 
  \textit{Austrian Seismic Network} [Data set]. 
  GeoSphere Austria / International Federation of Digital Seismograph Networks. 
  \url{https://doi.org/10.7914/SN/OE}

  \item \textbf{[OX] Istituto Nazionale di Oceanografia e di Geofisica Sperimentale -- OGS (2016)}. 
  \textit{North-East Italy Seismic Network} [Data set]. 
  International Federation of Digital Seismograph Networks. 
  \url{https://doi.org/10.7914/SN/OX}

  \item \textbf{[RF] University of Trieste (1993)}. 
  \textit{Friuli Venezia Giulia Accelerometric Network (RAF)} [Data set]. 
  International Federation of Digital Seismograph Networks. 
  \url{https://doi.org/10.7914/SN/RF}

  \item \textbf{[S] Swiss Seismological Service (SED) at ETH Z\"urich (2008)}. 
  \textit{Seismology at School Program, ETH Z\"urich (Seismo-at-school)} [Data set]. 
  ETH Z\"urich. 
  \url{https://doi.org/10.12686/SED/NETWORKS/S}

  \item \textbf{[SI] ZAMG and Autonomous Province of Bozen/Bolzano (2006)}. 
  \textit{Province S\"udtirol Network} [Data set]. 
  International Federation of Digital Seismograph Networks. 
  \url{https://www.fdsn.org/networks/detail/SI/}

  \item \textbf{[SL] Slovenian Environment Agency (ARSO) (1990)}. 
  \textit{Seismic Network of the Republic of Slovenia} [Data set]. 
  International Federation of Digital Seismograph Networks. 
  \url{https://doi.org/10.7914/SN/SL}

  \item \textbf{[ST] Geological Survey -- Provincia Autonoma di Trento (1981)}. 
  \textit{Trentino Seismic Network} [Data set]. 
  International Federation of Digital Seismograph Networks. 
  \url{https://doi.org/10.7914/SN/ST}

  \item \textbf{[ZO] Massa, M., Rizzo, A.~L., Lorenzetti, A., Lovati, S., D'Alema, E., Puglia, R., Carannante, S., \& Luzi, L. (2021)}. 
  \textit{Rete di monitoraggio multiparametrico del Garda -- PDnet (Version 1.0)} [Data set]. 
  Istituto Nazionale di Geofisica e Vulcanologia (INGV). 
  \url{https://doi.org/10.13127/SD/YHCFOMCBO_}
\end{enumerate}

\section*{Benchmark Machine-Learning Training Corpora}

The deep-learning phase pickers evaluated in this thesis were accessed through the standardized SeisBench repository (\cite{SeisBench}, \url{https://github.com/seisbench/seisbench}):
\begin{itemize}
  \item \textbf{INSTANCE (Italy)}: Michelini, A., Cianetti, S., Gaviano, S., Giunchi, C., Jozinovi\'c, D., \& Lauciani, V. (2021). \textit{INSTANCE -- the Italian seismic dataset for machine learning}. \textit{Earth System Science Data}, 13(12), 5509--5544. \url{https://doi.org/10.5194/essd-13-5509-2021}
  \item \textbf{STEAD (Global)}: Mousavi, S.~M., Sheng, Y., Zhu, W., \& Beroza, G.~C. (2019). \textit{STanford EArthquake Dataset (STEAD): A Global Data Set of Seismic Signals for AI}. \textit{IEEE Access}, 7, 179464--179476. \url{https://doi.org/10.1109/ACCESS.2019.2947848}
  \item \textbf{SCEDC (Southern California)}: Southern California Earthquake Center (2013). \textit{Southern California Seismic Network (SCSN) Waveform Data}. Caltech/USGS. \url{https://doi.org/10.7909/C3WD3xH1}
  \item \textbf{PhaseNet Original Weights}: Zhu, W., \& Beroza, G.~C. (2018). \textit{PhaseNet: a deep-neural-network-based seismic arrival-time picking method}. \textit{Geophysical Journal International}, 216(1), 261--273. \url{https://doi.org/10.1093/gji/ggy423}
  \item \textbf{Earthquake Transformer Original Weights}: Mousavi, S.~M., Ellsworth, W.~L., Zhu, W., Chuang, L.~Y., \& Beroza, G.~C. (2020). \textit{Earthquake transformer---an attentive deep-learning model for simultaneous earthquake detection and phase picking}. \textit{Nature Communications}, 11(1), 3952. \url{https://doi.org/10.1038/s41467-020-17591-w}
\end{itemize}

\section*{Code and Software Availability}

The complete computational software developed in this dissertation is open-source and structured across modular packages:
\begin{itemize}
  \item \textbf{OGS Pipeline Framework}: The Python source code (\texttt{OGS/src/}) implementing multi-format catalog parsing, FDSN downloading, bipartite graph matching (\ac{BPGMA}), local magnitude calculation, and plotting is maintained in the institutional repository:
  \url{https://github.com/kentanaka3/PhD}.
  \item \textbf{CINECA Leonardo HPC Orchestration}: All Slurm submission scripts, Hydra configuration templates, Makefile dependency graphs, and environment activation wrappers (\texttt{OGS/utils/Leonardo/}) are preserved within the repository under release tag \texttt{v1.0-UNITS-PhD}.
  \item \textbf{SeisBench Platform}: Pre-trained models and dataset interfaces were executed using SeisBench version 0.3.4 (\cite{SeisBench}).
  \item \textbf{Associator Codes}:
  \begin{itemize}
    \item \textbf{GaMMA}: Available at \url{https://github.com/AI4EPS/GaMMA} (\cite{GaMMA}).
    \item \textbf{PyOcto}: Available at \url{https://github.com/austin-bell/pyocto} (\cite{PyOcto}).
  \end{itemize}
  \item \textbf{Relocation Engine}: The NonLinLoc software suite (version 7.00) developed by Anthony Lomax is accessible at \url{https://alomax.free.fr/nlloc/} (\cite{NonLinLoc}).
\end{itemize}

